% Part: normal-modal-logic
% Chapter: axioms-systems
% Section: introduction

\documentclass[../../../include/open-logic-section]{subfiles}

\begin{document}

\olfileid[fr]{nml}{prf}{int}

\olsection{Introduction}

Nous disposons d'une sémantique du langage modal de base fondée sur
les modèles modaux. Une !!{formula} est \emph{valide} si elle est vraie
en tout monde de tout modèle; elle est valide relativement à une
classe de modèles ou de cadres si elle est vraie en tout monde de tout
modèle de la classe ou fondé sur un cadre de celle-ci. En logique, on
relie généralement ces caractérisations sémantiques de la validité à
une notion, issue de la théorie de la démonstration, de
!!{derivability}. Notre but est de définir, dans un système, une
notion de !!{derivability} telle qu'une !!{formula} soit
!!{derivable} si et seulement si elle est valide.

Les systèmes de !!{derivation} les plus simples et les plus anciens
sont les systèmes axiomatiques, dits \emph{de type Hilbert}. Pour de
nombreuses logiques modales, ils sont relativement faciles à
construire et simples à étudier sur le plan métathéorique, par
exemple pour démontrer leur correction et leur complétude. En
revanche, ils sont bien moins commodes que les systèmes de déduction
naturelle pour démontrer des !!{formula}s.

Dans un système de !!{derivation} de type Hilbert, une
!!{derivation} d'une !!{formula} est une suite de !!{formula}s qui
conduit de certains axiomes à cette formule au moyen de quelques
règles d'inférence. Pour que le système corresponde à la sémantique,
les formules !!{derivable}s doivent être vraies dans tous les modèles,
ou dans tous les modèles où les axiomes sont vrais. Nous isolerons
d'abord les propriétés requises des logiques modales dites
\emph{normales}. Pour ces logiques, deux règles d'inférence suffisent :
le modus ponens et la nécessitation. Les axiomes comprennent toutes
les instances par substitution des tautologies et, selon la logique
modale considérée, certains axiomes modaux. Même si nous ne retenons
que la classe de tous les modèles, il faut aussi compter comme axiomes
toutes les instances par substitution de~$\Ax{K}$\iftag{prvDiamond}{ et de~$\Dual$}{}.
On obtient ainsi la plus petite logique modale normale,~$\Log K$.

\begin{defn}
La règle du \emph{modus ponens} est le schéma d'inférence
\begin{prooftree}
\AxiomC{$!A$}
\AxiomC{$!A \lif !B$}
\RightLabel{\MP}
\BinaryInfC{$!B$}
\end{prooftree}
Nous disons que la !!{formula}~$!B$ \emph{découle de}~$!A$ et $!C$
par modus ponens si et seulement si $!C \ident !A \lif !B$.
\end{defn}

\begin{defn}
La règle de \emph{nécessitation} est le schéma d'inférence
\begin{prooftree}
\AxiomC{$!A$}
\RightLabel{\Nec}
\UnaryInfC{$\Box !A$}
\end{prooftree}
Nous disons que la !!{formula}~$!B$ découle de la !!{formula}~$!A$
par nécessitation si et seulement si $!B \ident \Box !A$.
\end{defn}

\begin{defn}
Une \emph{!!{derivation}} à partir d'un ensemble d'axiomes~$\Sigma$
est une suite de !!{formula}s $!B_1$, $!B_2$, \dots, $!B_n$ où chaque
$!B_i$ est :
\begin{enumerate}
\item une instance par substitution d'une tautologie;
\item une instance par substitution d'une !!{formula} de~$\Sigma$;
\item une formule qui découle de deux !!{formula}s $!B_j$, $!B_k$
  avec $j$, $k < i$ par modus ponens;
\item une formule qui découle d'une !!{formula}~$!B_j$ avec $j < i$
  par nécessitation.
\end{enumerate}
S'il existe une telle !!{derivation} avec $!B_n \ident !A$, nous
disons que $!A$ est \emph{!!{derivable} de $\Sigma$}, ce que nous
notons $\Sigma \Proves !A$.
\end{defn}

Lorsque $\Sigma$ contient~$\Ax{K}$\iftag{prvDiamond}{ et~$\Dual$}{},
les formules !!{derivable}s de~$\Sigma$ forment une logique modale
normale. Plus généralement, soit un modèle dans lequel toutes les
instances de substitution des axiomes de~$\Sigma$ sont vraies en
chaque monde. Toute !!{formula} !!{derivable} de~$\Sigma$ est alors
vraie dans ce modèle. C'est la \emph{correction} de la dérivation.
La \emph{complétude} consiste à établir la réciproque pour le
système considéré.
\end{document}
